\begin{textsample}{POS Dim 1 – rj – Score 32.00 – rj/rj\_inf\_e\_ef\_intro\_7.txt}  \label{ex:f1_pos_277}
1.6 Avaliação Na linguagem cotidiana, o \textbf{verbo} avaliar significa estimar, apreciar, examinar, \textbf{implicando} em coleta de informações sobre um determinado objeto e atribuição de valor ou qualidade ao mesmo.
Nesse processo, realizamos uma comparação entre o objeto e um determinado padrão previamente estabelecido como \textbf{parâmetro}, formulando um juízo de valor.
Na vida escolar, tradicionalmente, o resultado da avaliação tem servido para criar pequenas hierarquias.
Os estudantes são comparados e depois classificados em função de um padrão considerado de excelência.
A proximidade ou mesmo o distanciamento com relação a esse \textbf{parâmetro} define o êxito ou o fracasso.
Avaliar, entretanto, é mais do que isto, mais do que classificar os estudantes em aprovados ou reprovados, em bons ou ruins.
Uma importante função do processo \textbf{avaliativo} é possibilitar o desenvolvimento da aprendizagem.
A coleta de dados sobre a situação do estudante nos oferece informações que podem fundamentar novas decisões sobre o processo pedagógico.
A avaliação da aprendizagem não se esgota em si mesma.
Avaliamos para intervir, para agir e corrigir os rumos do nosso trabalho.
Nesse sentido, se a avaliação contribuir para o desenvolvimento dos estudantes, pode -se \textbf{dizer} que essa se transforma em uma ferramenta pedagógica para melhorar a aprendizagem e a qualidade do ensino.
Assim, é preciso destacar que em decorrência das interações em \textbf{aula}, o professor é capaz de observar mais criteriosamente seus estudantes, buscar formas de gerir as aprendizagens e reorganizar o planejamento das ações pedagógicas.
A avaliação precisa estar a serviço do estudante e não da classificação.
Para ser educativa, a avaliação precisa estar voltada para a formação integral dos estudantes e não somente para a sua instrução.
Assim, é importante refletir também, sobre duas dimensões \textbf{avaliativas} ( técnica e ética ) que não se confundem, mas se complementam e devem estar sempre presentes no processo \textbf{avaliativo}.
Afinal, a avaliação não se encerra com a aplicação de um instrumento e com a \textbf{análise} dos resultados obtidos.
Avaliar \textbf{implica} em tomar decisões para o futuro, pois o processo é contínuo para o aperfeiçoamento do ensino e uma ferramenta para planejar/avaliar/replanejar e contribuir para a concretização dos direitos de aprendizagem indicados na BNCC. 1.7.
Estrutura do Documento de Orientação Curricular ( DOC-RJ ) A estrutura deste Documento de Orientação Curricular é composta por textos de apresentação sobre as etapas Educação Infantil e Ensino Fundamental.
Além disso, tais etapas trazem em suas propostas \textbf{quadros} explicativos com os objetivos de aprendizagem e desenvolvimento da Educação Infantil e \textbf{quadro} de habilidades para cada componente curricular do Ensino Fundamental, de forma semelhante à BNCC.
Na etapa Educação Infantil, são apresentados seis direitos de aprendizagem e desenvolvimento ( Conviver, Brincar, Participar, Explorar, Expressar e Conhecer -se ) além do estabelecimento de cinco campos de experiências ( O eu, o outro e o nós / Corpo, gestos e movimentos / Traços, sons, cores e formas / Escuta, fala, pensamento e imaginação / Espaços, tempos, quantidades, relações e transformações ).
No Ensino Fundamental, cada área de conhecimento apresenta \textbf{competências} específicas de área cujas \textbf{competências} \textbf{explicitam} como as dez \textbf{competências} gerais podem ser identificadas em cada área.
Ainda, cada componente curricular também apresenta suas respectivas \textbf{competências} específicas a serem desenvolvidas pelos estudantes ao longo de seu percurso formativo.
Cada componente curricular também traz um conjunto de habilidades que tem como objetivo expressar as aprendizagens essenciais do processo de \textbf{ensino-aprendizagem} nas quais se encontram relacionados a objetos de conhecimento e organizados por unidades temáticas.
Nos \textbf{quadros} que apresentam as unidades temáticas, os objetos de conhecimento e as habilidades definidas para cada ano ( ou bloco de anos ), cada habilidade é identificada por um código alfanumérico cuja composição é a seguinte : Ressalte -se que novas habilidades foram inseridas no qual aparecem identificadas com o prefixo RJ.
Também no DOC-RJ, conforme exposto na BNCC e por compreender a importância de se apontar como norte para o desenvolvimento e avanço progressivo no percurso formativo de nossos estudantes, mantem -se as 10 \textbf{competências} gerais ( conforme \textbf{quadro} abaixo ) a serem desenvolvidas de forma integrada aos componentes curriculares ao longo de toda a educação básica tendo em \textbf{vista} a aquisição de conhecimentos e habilidades e a formação de atitudes e valores com \textbf{vistas} à formação de estudantes verdadeiramente preparados para essenciais para a atuação no \textbf{século} XXI.
Enfatiza -se que nesta tessitura de conhecimento que as \textbf{competências} gerais definidas na BNCC como \textbf{mobilização} de conhecimentos, habilidades, atitudes e valores para resolver \textbf{demandas} \textbf{complexas} da vida cotidiana, do pleno exercício da cidadania e do mundo do trabalho, reafirmam o direito de aprender.
Por meio delas, a BNCC assim como o DOC-RJ institui seu compromisso com a educação integral, desenvolvendo a ideia de escola como espaço para a democracia inclusiva e respeitando a diferença e a diversidade.
As \textbf{competências} tem como \textbf{foco} o desenvolvimento : do conhecimento ; pensamento científico, \textbf{crítico} e criativo ; repertório cultural ; comunicação ; cultura \textbf{digital} ; trabalho e projeto de vida ; argumentação ; autoconhecimento autocuidado ; empatia e cooperação e responsabilidade e cidadania.
Dez \textbf{Competências} Gerais da Educação Básica 1.
Valorizar e utilizar os conhecimentos historicamente construídos sobre o mundo físico, social, cultural e \textbf{digital} para entender e explicar a realidade, continuar aprendendo e colaborar para a construção de uma sociedade justa, democrática e inclusiva. 2.
Exercitar a curiosidade intelectual e recorrer à \textbf{abordagem} própria das \textbf{ciências}, incluindo a investigação, a reflexão, a \textbf{análise} \textbf{crítica}, a imaginação e a criatividade, para investigar causas, elaborar e testar hipóteses, formular e resolver \textbf{problemas} e criar soluções ( inclusive tecnológicas ) com base nos conhecimentos das diferentes áreas. 3.
Valorizar e fruir as diversas manifestações artísticas e culturais, das locais às mundiais, e também participar de práticas diversificadas da produção \textbf{artístico-cultural}. 4.
Utilizar diferentes linguagens – verbal ( oral ou \textbf{visual-motora}, como \textbf{Libras}, e escrita ), corporal, visual, sonora e \textbf{digital} –, bem como conhecimentos das linguagens artística, matemática e científica, para se expressar e partilhar informações, experiências, ideias e sentimentos em diferentes contextos e produzir sentidos que levem ao entendimento mútuo. 5.
Compreender, utilizar e criar \textbf{tecnologias} \textbf{digitais} de informação e comunicação de forma \textbf{crítica}, significativa, \textbf{reflexiva} e ética nas diversas práticas sociais ( incluindo as escolares ) para se comunicar, acessar e disseminar informações, produzir conhecimentos, resolver \textbf{problemas} e exercer protagonismo e \textbf{autoria} na vida pessoal e coletiva. 6.
Valorizar a diversidade de saberes e vivências culturais e apropriar -se de conhecimentos e experiências que lhe possibilitem entender as relações próprias do mundo do trabalho e fazer escolhas alinhadas ao exercício da cidadania e ao seu projeto de vida, com liberdade, autonomia, consciência \textbf{crítica} e responsabilidade. 7. \textbf{Argumentar} com base em fatos, dados e informações confiáveis, para formular, negociar e defender ideias, pontos de \textbf{vista} e decisões comuns que respeitem e promovam os direitos humanos, a consciência \textbf{socioambiental} e o consumo responsável em âmbito local, regional e global, com \textbf{posicionamento} ético em relação ao cuidado de si mesmo, dos outros e do planeta. 8.
Conhecer -se, apreciar -se e cuidar de sua saúde física e emocional, compreendendo -se na diversidade humana e reconhecendo suas emoções e as dos outros, com autocrítica e \textbf{capacidade} para lidar com elas. 9.
Exercitar a empatia, o diálogo, a resolução de conflitos e a cooperação, fazendo -se respeitar e promovendo o respeito ao outro e aos direitos humanos, com acolhimento e valorização da diversidade de indivíduos e de grupos sociais, seus saberes, identidades, culturas e potencialidades, sem preconceitos de qualquer natureza. 10.
Agir pessoal e coletivamente com autonomia, responsabilidade, flexibilidade, resiliência e determinação, tomando decisões com base em princípios éticos, democráticos, inclusivos, sustentáveis e solidários.

% matched lemmas: abordagem, análise, argumentar, artístico-cultural, aula, autoria, avaliativo, capacidade, ciência, competência, complexo, crítico, demanda, digital, dizer, ensino-aprendizagem, explicitar, foco, implicar, libra, mobilização, parâmetro, posicionamento, problema, quadro, reflexivo, socioambiental, século, tecnologia, verbo, vista, visual-motor
\end{textsample}
